\section{Kausale Überprüfung mit Activation Patching}

Alle bisherigen Verfahren liefern Korrelationen: Sie zeigen, dass eine Information vorhanden ist oder genutzt zu werden scheint, aber nicht, dass eine bestimmte Komponente für die Vorhersage tatsächlich notwendig ist. Diesen Nachweis liefert das \textbf{Activation Patching}, das in der mechanistischen Interpretierbarkeit als kausaler Goldstandard gilt \cite{meng2022locating, goldowskydill2023localizing}. Das Prinzip ist ein kontrolliertes Experiment: Ein \emph{sauberer} und ein \emph{korrumpierter} Durchlauf unterscheiden sich nur im sentimenttragenden Wort. Im korrumpierten Durchlauf wird dann gezielt die \gls{activation} einer einzelnen Schicht an einer einzelnen Token-Position durch die zuvor gespeicherte saubere Version ersetzt. Verändert sich die Vorhersage dadurch wieder in Richtung des sauberen Ergebnisses, enthält genau diese Schicht und Position die entscheidende Information. Normalisiert bedeutet ein Effekt von 0 keine Wiederherstellung, ein Effekt von 1 eine vollständige Wiederherstellung der sauberen Vorhersage. Technisch läuft das über PyTorch Forward Hooks: Zunächst werden alle sauberen \glspl{activation} zwischengespeichert, danach wird systematisch über jede Kombination aus Schicht und Position gepatcht und der jeweilige Effekt gemessen, und zuletzt werden die Ergebnisse als Heatmap dargestellt und die einflussreichsten Punkte herausgefiltert.

Als Prompt-Paar dient erneut das bereits aus den vorigen Kapiteln bekannte Beispiel \texttt{The food was delicious/disgusting and I felt very} mit Zielwort „satisfied“, derselben Formulierung, an der schon der Logit Lens und die Attention-Analyse untersucht wurden, sodass sich die drei Methoden direkt miteinander vergleichen lassen. Im sauberen Durchlauf erreicht dieses Wort einen Logit von 15,78 (Wahrscheinlichkeit 10,2~Prozent), im korrumpierten Durchlauf nur 9,21 (Wahrscheinlichkeit unter 0,1~Prozent). Aus der Erwartung der vorigen Kapitel heraus sollte der stärkste kausale Effekt in den oberen Schichten 17 bis 23 liegen, dort, wo Logit Lens und Attention die auffälligsten Muster zeigten.

\begin{figure}[htbp]
  \centering
  \includegraphics[width=\textwidth]{activation_patching_heatmap.png}
  \caption{Normalisierter Activation-Patching-Effekt über alle Schichten und Token-Positionen.}
  \label{fig:patch}
\end{figure}

\begin{table}[htbp]
  \centering
  \caption{Die fünf einflussreichsten Patching-Punkte.}
  \label{tab:patchtop}
  \begin{tabular}{lrrlr}
    \toprule
    Rang & Schicht & Position & Token & Effekt \\
    \midrule
    1 & 2 & 3 & delicious & 1,023 \\
    2 & 1 & 3 & delicious & 1,021 \\
    3 & 0 & 3 & delicious & 1,017 \\
    4 & 3 & 3 & delicious & 1,016 \\
    5 & 6 & 3 & delicious & 1,016 \\
    \bottomrule
  \end{tabular}
\end{table}

Das tatsächliche Ergebnis überrascht: Neun der zehn stärksten Punkte liegen exakt an Position 3, dem Wort „delicious“, verteilt über die frühen Schichten 0 bis 8 (Tabelle~\ref{tab:patchtop}, Abbildung~\ref{fig:patch}). Die kausal entscheidende Information sitzt also nicht in den oberen, sondern bereits in den allerfrühesten Schichten an der Token-Repräsentation selbst, passend zur linearen Trennbarkeit der Embeddings aus Kapitel 4, aber im Widerspruch zur Erwartung aus Logit Lens und Attention.

Ergänzend wird auf Ebene einzelner Köpfe gepatcht: Statt der gesamten Schicht-Aktivierung wird nur der Anteil eines einzelnen Attention-Kopfes an Position 3 ersetzt (Abbildung~\ref{fig:patchhead}). Der stärkste Einzelkopf liegt in Schicht 5, Kopf 12, mit einem Effekt von lediglich 0,019, rund zwei Größenordnungen schwächer als der Effekt auf Schichtebene. Kein einzelner Kopf erklärt den beobachteten Effekt auch nur annähernd, und insbesondere taucht der zuvor auffällige Kopf aus Kapitel 6 (Schicht 6, Kopf 3) hier nicht auf.

\begin{figure}[htbp]
  \centering
  \includegraphics[width=0.7\textwidth]{activation_patching_head_heatmap.png}
  \caption{Patching auf Ebene einzelner Attention Heads an Position 3 (Token „delicious“).}
  \label{fig:patchhead}
\end{figure}

Für dieses eine Prompt-Paar bestätigt sich damit der kausale Teil von Hypothese H1 vollständig: Ein Eingriff am Stimmungswort stellt die Vorhersage wieder her. Hypothese H2 bestätigt sich nur teilweise: Die Position ist scharf lokalisiert, die Schicht dagegen nicht, sondern verteilt sich über ein Band aus neun frühen Schichten statt auf einen einzelnen Punkt. Da nur ein Beispiel geprüft wurde, ist offen, ob sich das Muster über den gesamten CAD-Datensatz verallgemeinern lässt.
